%% kapa: exam=1 round=21 year=2010 period=1 form=A subject=civil_law subject_ko=민법 seq=32 number=32
%% source: Q-Net 감정평가사 기출문제 (cst00309) artlSeq=5024181 file=extracted/1008/5024181/2010년 제21회 감정평가사 1차시험 기출문제/1교시.hwp
%% conversion: kordoc 4.24.51 -> tools/gen1.py (auto); review=proofread-sonnet; flags=none
\begin{question}{32}
민법상 법인에 관한 설명으로 옳지 않은 것은? (다툼이 있으면 판례에 의함)
\begin{choices}
  \item 재단법인의 기본재산의 처분은 정관변경을 요하는 것이므로, 주무관청의 허가가 없으면 그 처분행위는 무효이다.
  \item 사단법인의 사원들이 사원총회의 결의로 정관의 규범적인 의미 내용과 다른 해석을 한 경우, 그 해석은 사단법인의 사원들이나 법원을 구속하는 효력이 없다.
  \item 법률이 정하는 일정한 조직과 요건을 갖추었을 경우 주무관청은 반드시 그 법인의 설립을 허가하여야 한다.
  \item 법인 아닌 사단의 대표자가 직무에 관하여 타인에게 손해를 가한 때에도 법인의 불법행위책임에 관한 규정이 유추 적용된다.
  \item 감사는 법인의 재산상황이나 이사의 업무집행을 감사하는 직무 등을 하며, 필요시 사원총회를 소집할 권리가 있다.
\end{choices}
\end{question}
\answer{3}
